% Осень26!B4; нумерация внутри столбца без пустых строк.
% Источник: лекция 4, §15 (строка 504)
\definitionticket{02}{Положительно-, отрицательно-, незнако- определенная квадратичная форма}
\label{definition:02}

\begin{definition}
\textbf{Квадратичная форма} на $\mathbb R^m$ --- функция
\[
  Q(h)=\sum_{i,j=1}^{m}a_{ij}h_i h_j,
  \qquad a_{ij}\in\mathbb R.
\]
Она называется:
\begin{itemize}
  \item \textbf{положительно определённой}, если $Q(h)>0$ для всех $h\ne0$;
  \item \textbf{отрицательно определённой}, если $Q(h)<0$ для всех $h\ne0$;
  \item \textbf{незнакоопределённой}, если существуют $h_1,h_2$ такие, что $Q(h_1)>0$ и $Q(h_2)<0$.
\end{itemize}
\end{definition}
